\chapter{Maxwell--Boltzmann Distribution}

\section*{Introduction}

The Maxwell--Boltzmann distribution describes the statistical distribution of speeds among particles in an ideal gas at thermal equilibrium. Developed by James Clerk Maxwell (1860) and refined by Ludwig Boltzmann (1868), it is the cornerstone of the kinetic theory of gases and one of the earliest triumphs of statistical mechanics. The distribution emerges from the assumption that each particle's kinetic energy is distributed according to the Boltzmann factor $e^{-E/kT}$, combined with the three-dimensional nature of velocity space.

In an ideal gas at temperature $T$, particles move randomly and collide elastically. At equilibrium, the number of particles with speeds between $v$ and $v + dv$ follows the Maxwell--Boltzmann distribution $f(v)\,dv$. The distribution is asymmetric: it rises from zero at $v = 0$, peaks at the most probable speed $v_{\text{mp}} = \sqrt{2kT/m}$, and has a long tail toward high speeds. The three characteristic speeds---most probable, mean, and root-mean-square---are all different and increase with temperature, reflecting the thermal agitation of the gas.
\section*{Fundamental Equation Used}

\begin{equation}
\boxed{\; f(v) = 4\pi n \left(\frac{m}{2\pi kT}\right)^{3/2} v^{2}\,e^{-mv^{2}/2kT}\;}
\label{eq:maxwell-boltzmann}
\end{equation}

The three characteristic speeds are:
\begin{align}
v_{\text{mp}} &= \sqrt{\frac{2kT}{m}}, \qquad \bar{v} = \sqrt{\frac{8kT}{\pi m}}, \qquad v_{\text{rms}} = \sqrt{\frac{3kT}{m}}.
\label{eq:mb-speeds}
\end{align}
\section*{Assumptions}

\textbf{Assumptions:} The gas is ideal (non-interacting point particles); the system is in thermal equilibrium; the number of particles is large.


\textbf{Limitations:} Does not apply to quantum gases (Bose--Einstein or Fermi--Dirac) at low temperatures or high densities; does not account for internal molecular degrees of freedom.


\section*{Mathematical Derivation}

\textbf{Step 1: Start from the Boltzmann factor in velocity space.}

At thermal equilibrium, the probability that a particle has velocity $\mathbf{v} = (v_x, v_y, v_z)$ is proportional to the Boltzmann factor:
\begin{equation}
f(\mathbf{v})\,d^{3}v \propto e^{-mv^{2}/2kT}\,d^{3}v,
\end{equation}
where $v^{2} = v_x^{2} + v_y^{2} + v_z^{2}$. In Cartesian velocity-space volume element $d^{3}v = dv_x\,dv_y\,dv_z$.

\textbf{Step 2: Find the normalization constant.}

Requiring $\int f(\mathbf{v})\,d^{3}v = n$ (the total number density):
\begin{equation}
n = A\int_{-\infty}^{\infty}\!\!\int_{-\infty}^{\infty}\!\!\int_{-\infty}^{\infty} e^{-m(v_x^{2}+v_y^{2}+v_z^{2})/2kT}\,dv_x\,dv_y\,dv_z.
\end{equation}
The integral factors into three identical Gaussian integrals:
\begin{equation}
\int_{-\infty}^{\infty} e^{-mv_x^{2}/2kT}\,dv_x = \sqrt{\frac{2\pi kT}{m}},
\end{equation}
so the normalization gives $n = A\,(2\pi kT/m)^{3/2}$, hence $A = n\,(m/2\pi kT)^{3/2}$.

\textbf{Step 3: Transform to spherical velocity coordinates.}

The velocity-space distribution is isotropic, so we switch to spherical coordinates $(v, \theta, \phi)$ where $d^{3}v = v^{2}\sin\theta\,dv\,d\theta\,d\phi$. Integrating over the angular variables:
\begin{equation}
\int_{0}^{2\pi}\!\!\int_{0}^{\pi} \sin\theta\,d\theta\,d\phi = 4\pi.
\end{equation}
The speed distribution is then:
\begin{equation}
f(v)\,dv = 4\pi n\left(\frac{m}{2\pi kT}\right)^{3/2} v^{2}\,e^{-mv^{2}/2kT}\,dv.
\end{equation}

\textbf{Step 4: Find the most probable speed by differentiation.}

Setting $df/dv = 0$:
\begin{equation}
\frac{df}{dv} \propto \frac{d}{dv}\!\left(v^{2}e^{-mv^{2}/2kT}\right) = 2v\,e^{-mv^{2}/2kT} - \frac{mv^{3}}{kT}\,e^{-mv^{2}/2kT} = 0.
\end{equation}
Factoring $ve^{-mv^{2}/2kT}$:
\begin{equation}
2 - \frac{mv^{2}}{kT} = 0 \quad\Longrightarrow\quad v_{\text{mp}} = \sqrt{\frac{2kT}{m}}.
\end{equation}

\textbf{Step 5: Compute the mean speed.}

\begin{equation}
\bar{v} = \frac{1}{n}\int_{0}^{\infty} v\,f(v)\,dv = 4\pi\left(\frac{m}{2\pi kT}\right)^{3/2}\int_{0}^{\infty} v^{3}\,e^{-mv^{2}/2kT}\,dv.
\end{equation}
Using the integral $\int_{0}^{\infty} v^{3}e^{-av^{2}}\,dv = \frac{1}{2a^{2}}$ with $a = m/2kT$:
\begin{equation}
\bar{v} = 4\pi\left(\frac{m}{2\pi kT}\right)^{3/2}\cdot\frac{1}{2}\cdot\frac{(2kT)^{2}}{m^{2}} = \sqrt{\frac{8kT}{\pi m}}.
\end{equation}

\textbf{Step 6: Compute the root-mean-square speed.}

\begin{equation}
v_{\text{rms}}^{2} = \frac{1}{n}\int_{0}^{\infty} v^{2}\,f(v)\,dv = 4\pi\left(\frac{m}{2\pi kT}\right)^{3/2}\cdot\frac{(2kT)^{5/2}}{4m^{5/2}\cdot\Gamma(5/2)} = \frac{3kT}{m},
\end{equation}
giving $v_{\text{rms}} = \sqrt{3kT/m}$. The three characteristic speeds satisfy $v_{\text{mp}} < \bar{v} < v_{\text{rms}}$.

\section*{Characteristics of the Equation}

\begin{itemize}
  \item \textbf{Normalization:} $\int_{0}^{\infty} f(v)\,dv = n$, the total number density of particles.
  \item \textbf{Temperature dependence:} As $T$ increases, the distribution broadens and the peak shifts to higher $v$; the area under the curve is conserved ($n$ fixed).
  \item \textbf{Mass dependence:} Heavier molecules have narrower distributions at the same $T$; lighter molecules have broader distributions.
  \item \textbf{Asymmetry:} The distribution has a positive skew: $\bar{v} > v_{\text{mp}}$ and $v_{\text{rms}} > \bar{v}$.
\end{itemize}
\section*{Solution Methods}

\begin{enumerate}
  \item \textbf{From the Boltzmann factor:} The probability of a state with energy $E = \frac{1}{2}mv^{2}$ is proportional to $e^{-E/kT}$. In three dimensions, the velocity-space volume element is $4\pi v^{2}\,dv$, giving $f(v) \propto v^{2}\,e^{-mv^{2}/2kT}$. The normalization constant is found by integration.
  \item \textbf{From the partition function:} The single-particle partition function for translational motion is $Z_{1} = (2\pi mkT/h^{2})^{3/2}V$. The speed distribution follows from $f(v) = 4\pi v^{2}(m/2\pi kT)^{3/2}e^{-mv^{2}/2kT}$.
  \item \textbf{From the entropy maximum:} Maximizing the entropy $S = -k\sum p_{i}\ln p_{i}$ subject to constraints on particle number and total energy yields the Boltzmann distribution; restricting to translational kinetic energy gives the Maxwell--Boltzmann distribution.
\end{enumerate}
\section*{Verifications}

The Maxwell--Boltzmann distribution has been verified through molecular beam experiments, Doppler broadening measurements, and neutron scattering. The distribution of speeds in a gas can be directly measured using time-of-flight techniques or laser Doppler velocimetry, and results agree with theory to high precision for gases at moderate temperatures and pressures.
\section*{Applications}

\begin{itemize}
  \item \textbf{Atmospheric escape:} Particles in the tail of the distribution with $v$ exceeding the escape velocity leave a planet's atmosphere. This explains why light gases (H$_{2}$, He) are rare on Earth but abundant on gas giants.
  \item \textbf{Molecular beams:} Effusion through a small hole selects particles with a flux-weighted distribution $f_{\text{flux}}(v) \propto v^{3}\,e^{-mv^{2}/2kT}$, which is shifted toward higher speeds.
  \item \textbf{Chemical kinetics:} The Arrhenius activation energy barrier is overcome by particles in the high-energy tail of the distribution.
  \item \textbf{Spectroscopy:} Doppler broadening of spectral lines directly reflects the Maxwell--Boltzmann speed distribution of emitting or absorbing atoms.
\end{itemize}
\section*{What Richard Feynman Would Have Asked}

\subsection*{Plain-English Translation}
\textit{``What is the Maxwell--Boltzmann distribution really telling us about the world?''}

In a gas, molecules are bouncing around chaotically, colliding billions of times per second. The distribution tells you how the speeds are spread out at equilibrium: most molecules have moderate speeds, a few are crawling, and a few are zooming very fast. The shape of this spread is a universal consequence of three things: energy is quadratic in velocity ($E = \frac{1}{2}mv^{2}$), the probability of a state decreases exponentially with energy ($e^{-E/kT}$), and we live in three dimensions (which gives the $v^{2}$ factor). Change any one of these ingredients, and the distribution changes.

\subsection*{Localized Mechanics}
\textit{``How does the distribution emerge from the microscopic dynamics of particle collisions?''}

Imagine a gas of hard spheres undergoing elastic collisions. Each collision exchanges energy between two particles. Over time, the system explores the space of all possible energy distributions. The unique distribution that maximizes the number of microstates (entropy) consistent with a fixed total energy is the Boltzmann distribution $p(E) \propto e^{-E/kT}$. For free particles with $E = \frac{1}{2}mv^{2}$, this becomes the Maxwell--Boltzmann distribution in velocity space. The key insight is that any initial distribution will evolve toward this equilibrium form through collisions, because it is the maximum-entropy state. This is a statistical statement, not a dynamical one: it does not require that every collision obey the distribution, only that the aggregate effect of many collisions drives the system toward it.

\subsection*{Testing Extreme Limits}
\textit{``What happens to the distribution at extreme temperatures or for extreme particle masses?''}

At very low $T$, the de Broglie wavelength $\lambda_{\text{dB}} = h/\sqrt{2\pi mkT}$ becomes comparable to the inter-particle spacing, and the Maxwell--Boltzmann distribution must be replaced by Bose--Einstein or Fermi--Dirac statistics. For electrons in a metal at room temperature, $\lambda_{\text{dB}} \sim 4$~nm, much larger than the lattice spacing, so quantum statistics are essential. At very high $T$ (e.g., in the interior of stars), the distribution extends to relativistic speeds and must be modified. For extremely massive particles (e.g., dust grains in a gas), the distribution is extremely narrow---heavy particles barely move at room temperature. In the limit $m \to \infty$, the distribution becomes a delta function at $v = 0$.

\subsection*{Dimensionality and Geometry}
\textit{``How does the dimensionality of space affect the shape of the distribution?''}

In $d$-dimensional space, the velocity-space volume element is proportional to $v^{d-1}\,dv$ (the surface area of a $(d-1)$-sphere). The speed distribution becomes $f(v) \propto v^{d-1}\,e^{-mv^{2}/2kT}$. In one dimension, $f(v) \propto e^{-mv^{2}/2kT}$ (Gaussian, no $v^{2}$ factor), which peaks at $v = 0$. In three dimensions, $f(v) \propto v^{2}\,e^{-mv^{2}/2kT}$, which peaks at $v_{\text{mp}} = \sqrt{2kT/m}$. In higher dimensions, the peak shifts further outward. The physical reason is that there are more ways to have a large speed in higher dimensions (more directions), so the volume factor ``wins'' against the exponential suppression at moderate speeds. The three-dimensional case is the physically relevant one for our universe, but the mathematical generalization is illuminating.

\subsection*{Cross-Disciplinary Connection}
\textit{``Where does this same mathematical structure---a Gaussian multiplied by a power-law volume factor---appear in other fields?''}

The Maxwell--Boltzmann distribution is a chi-squared distribution with three degrees of freedom, a structure ubiquitous in statistics. In finance, the distribution of returns in the Black--Scholes model is log-normal, which is related through a variable transformation. In biology, the distribution of protein folding energies or DNA binding energies often follows a Boltzmann-like exponential. In machine learning, the softmax function $e^{-E_{i}/T}/\sum e^{-E_{j}/T}$ is a discrete Boltzmann distribution, and simulated annealing uses a ``temperature'' parameter to explore energy landscapes. The deep structure is that of exponential families in statistics, of which the Maxwell--Boltzmann distribution is the most physically transparent example.

% ============================================================
% CHAPTER 42 — Mirror Equation
% ============================================================
\section*{FAQ}

\begin{enumerate}
\item \textbf{Why does the distribution peak at a speed greater than zero, even though the most probable kinetic energy is zero?}
The speed distribution includes the velocity-space volume factor $4\pi v^{2}$, which vanishes at $v = 0$. Even though $e^{-mv^{2}/2kT}$ is largest at $v = 0$, the factor $v^{2}$ suppresses the probability at low speeds, pushing the peak to $v_{\text{mp}} > 0$.

\item \textbf{Does the distribution apply to liquids and solids?}
The translational Maxwell--Boltzmann distribution applies to the center-of-mass motion of atoms in any phase, but in liquids and solids the particles are not free, and intermolecular interactions must be included. The distribution is quantitatively accurate only for dilute gases.

\item \textbf{What happens at very high temperatures?}
At very high $T$, relativistic effects become important ($kT \sim mc^{2}$), and the non-relativistic Maxwell--Boltzmann distribution must be replaced by its relativistic generalization.
\end{enumerate}
\section*{Important Bibliography}

\begin{enumerate}
\item Pathria, R.K. and Beale, P.D., \textit{Statistical Mechanics}, 4th ed. (Academic Press, 2021), Chapter 2.
\item Reif, F., \textit{Fundamentals of Statistical and Thermal Physics} (McGraw-Hill, 1965), Chapter 6.
\item Tolman, R.C., \textit{The Principles of Statistical Mechanics} (Dover, 1987), Chapter 3.
\item Huang, K., \textit{Statistical Mechanics}, 2nd ed. (Wiley, 1987), Chapter 6.
\end{enumerate}


